\documentclass[14pt,dvipdfmx]{jsarticle}
\input{lecture-preamble.tex}

\begin{document}

\lectureheader{数学Ⅱ}{5}{指数関数と対数関数}{1}{指数関数}{2-C}{指数関数を含む方程式,不等式}
  {159}{指数関数を含む方程式，不等式について考えよう}

% 同じ節の続き（1-2-B までの番号を引き継ぐ）
\setcounter{definition}{3}
\setcounter{theorem}{2}
\setcounter{technique}{2}

% ============ 1ページ目（表・左）指数方程式・例題 ============
\vspace{-1.5zw}% 目標の枠と見出しの間を詰める
\heading{指数関数を含む方程式}
指数関数を含む方程式を解いてみよう．

そもそも「方程式を解く」とは，具体的にどういうことだろうか．

つまり「方程式を解く」ということは「\fitblankbf{x = \text{○}}」の形まで持っていくことである．

\begin{exampleprob}{3}
  次の方程式を解け。
  \begin{tasks}[label=(\arabic*), label-width=1.6zw, item-indent=2zw, after-item-skip=0.15em](2)
    \task $8^x = 4$
    \task $9^x = 3^{x+1}$
  \end{tasks}
\end{exampleprob}

\begin{answerbox}{(1)}[9]
  方程式を変形すると $2^{3x}=2^2$\\
  よって $3x=2$\\
  これを解いて
  \boxanswer{$x=\dfrac{2}{3}$}
\end{answerbox}
\begin{answerbox}{(2)}[8]
  方程式を変形すると $3^{2x}=3^{x+1}$\\
  よって $2x=x+1$\\
  これを解いて
  \boxanswer{$x=1$}
\end{answerbox}
\vspace{0.3em}

\begin{techniquebox}{指数を含む方程式}
  \begin{itemize}[leftmargin=1.5zw, itemsep=0.15em]
    \item 両辺の\fitblankbf{\text{底}}をそろえて，\fitblankbf{\text{指数}}を比べる．
  \end{itemize}
\end{techniquebox}

\vfill
\nextpage

% ============ 2ページ目（表・右）指数方程式・練習 ============
\begin{practiceprob}{11}
  次の方程式を解け。
  \begin{tasks}[label=(\arabic*), label-width=1.6zw, item-indent=2zw, after-item-skip=0.15em](3)
    \task $4^x = 8$
    \task $8^x = \dfrac{1}{16}$
    \task $27^x = 3^{2-x}$
  \end{tasks}
\end{practiceprob}

\begin{answerbox}{(1)}[7]
  方程式を変形すると $2^{2x}=2^3$\\
  よって $2x=3$
  \boxanswer{$x=\dfrac{3}{2}$}
\end{answerbox}
\begin{answerbox}{(2)}[7]
  方程式を変形すると $2^{3x}=2^{-4}$\\
  よって $3x=-4$
  \boxanswer{$x=-\dfrac{4}{3}$}
\end{answerbox}
\begin{answerbox}{(3)}[7]
  方程式を変形すると $3^{3x}=3^{2-x}$\\
  よって $3x=2-x$
  \boxanswer{$x=\dfrac{1}{2}$}
\end{answerbox}
\vspace{0.5em}

\heading{演習問題 指数方程式}
\begin{exerciseprob}{1}
  次の方程式を解け。
  \begin{tasks}[label=(\arabic*), label-width=1.6zw, item-indent=2zw, after-item-skip=0.15em](3)
    \task $2^x = 256$
    \task $\left(\dfrac{1}{3}\right)^x = 81$
    \task $(\sqrt{2})^{2x-1} = 0.25^x$
  \end{tasks}
\end{exerciseprob}

\begin{answerbox}{(1)}[6]
  $256=2^8$ だから $2^x=2^8$ \\
  \boxanswer{$x=8$}
\end{answerbox}
\begin{answerbox}{(2)}[6]
  $\left(\dfrac{1}{3}\right)^x=3^{-x}$，$81=3^4$ だから $3^{-x}=3^4$ \\
  よって $-x=4$
  \boxanswer{$x=-4$}
\end{answerbox}
\begin{answerbox}{(3)}[6]
  $(\sqrt{2})^{2x-1}=2^{x-\frac{1}{2}}$，$0.25^x=2^{-2x}$ だから $x-\dfrac{1}{2}=-2x$ \\
  \boxanswer{$x=\dfrac{1}{6}$}
\end{answerbox}

\vfill
\nextpage

% ============ 3ページ目（裏・左）指数不等式・例題 ============
\heading{指数関数を含む不等式}
指数関数を含む不等式を解いてみよう．基本的な考え方は指数方程式を解くときと同じである．

\begin{exampleprob}{4}
  次の不等式を解け。
  \begin{tasks}[label=(\arabic*), label-width=1.6zw, item-indent=2zw, after-item-skip=0.15em](2)
    \task $2^x \geqq 8$
    \task $\left(\dfrac{1}{3}\right)^{x+1} < \left(\dfrac{1}{9}\right)^x$
  \end{tasks}
\end{exampleprob}

\begin{answerbox}{(1)}[6]
  不等式を変形すると $2^x\geqq 2^3$\\
  底 $2$ は $1$ よりも大きいので
  \boxanswer{$x\geqq 3$}
\end{answerbox}
\begin{answerbox}{(2)}[8]
  不等式を変形すると $\left(\dfrac{1}{3}\right)^{x+1}<\left(\dfrac{1}{3}\right)^{2x}$\\
  底 $\dfrac{1}{3}$ は $1$ よりも小さいので
  \boxanswer{$x<1$}
\end{answerbox}

\begin{practiceprob}{12}
  次の不等式を解け。
  \begin{tasks}[label=(\arabic*), label-width=1.6zw, item-indent=2zw, after-item-skip=0.15em](3)
    \task $3^x < 81$
    \task $\left(\dfrac{1}{2}\right)^x \geqq \dfrac{1}{32}$
    \task $2^{3x-4} > \left(\dfrac{1}{4}\right)^x$
  \end{tasks}
\end{practiceprob}

\begin{answerbox}{(1)}[6]
  不等式を変形すると $3^x<3^4$\\
  底 $3$ は $1$ よりも大きいので
  \boxanswer{$x<4$}
\end{answerbox}
\begin{answerbox}{(2)}[8]
  不等式を変形すると $\left(\dfrac{1}{2}\right)^x\geqq\left(\dfrac{1}{2}\right)^5$\\
  底 $\dfrac{1}{2}$ は $1$ よりも小さいので
  \boxanswer{$x\leqq 5$}
\end{answerbox}
\begin{answerbox}{(3)}[7]
  不等式を変形すると $2^{3x-4}>2^{-2x}$\\
  底 $2$ は $1$ よりも大きいので
  \boxanswer{$x>\dfrac{4}{5}$}
\end{answerbox}



\vfill
\nextpage

% ============ 4ページ目（裏・右）まとめ，演習問題 ============
\begin{techniquebox}{指数を含む不等式}
  \begin{itemize}[leftmargin=1.5zw, itemsep=0.15em]
    \item 両辺の\fitblankbf{\text{底}}をそろえて，\fitblankbf{\text{指数}}を比べる．
    \item 不等式では，底が $0<a<1$ のとき不等号の向きが\fitblankbf{\text{逆}}になる．
  \end{itemize}
\end{techniquebox}

\begin{summarybox}
  \begin{enumerate}[label=\textbf{\arabic*}, leftmargin=1.5zw, itemsep=0.2em]
    \item 指数を含む方程式・不等式は，\fitblankbf{\text{底}}をそろえて\fitblankbf{\text{指数}}を比べよう．
    \item 不等式では，底が $0<a<1$ のとき不等号の向きが\fitblankbf{\text{逆}}になることに注意しよう．
  \end{enumerate}
\end{summarybox}

\heading{演習問題 指数不等式}
\begin{exerciseprob}{2}
  次の不等式を解け。
  \begin{tasks}[label=(\arabic*), label-width=1.6zw, item-indent=2zw, after-item-skip=0.15em](3)
    \task $5^x \leqq \dfrac{1}{125}$
    \task $\left(\dfrac{1}{2}\right)^{x+1} > \dfrac{1}{4}$
    \task $(0.2)^{2x} > 5\sqrt{5}$
  \end{tasks}
\end{exerciseprob}

\begin{answerbox}{(1)}[6]
  $\dfrac{1}{125}=5^{-3}$ だから $5^x\leqq 5^{-3}$ \\
  底 $5>1$ なので向きはそのまま
  \boxanswer{$x\leqq -3$}
\end{answerbox}
\begin{answerbox}{(2)}[8]
  $\dfrac{1}{4}=\left(\dfrac{1}{2}\right)^2$ だから $\left(\dfrac{1}{2}\right)^{x+1}>\left(\dfrac{1}{2}\right)^2$ \\
  底 $0<\dfrac{1}{2}<1$ だから $x+1<2$
  \boxanswer{$x<1$}
\end{answerbox}
\begin{answerbox}{(3)}[8]
  $(0.2)^{2x}=5^{-2x}$，$5\sqrt{5}=5^{\frac{3}{2}}$ だから $5^{-2x}>5^{\frac{3}{2}}$ \\
  底 $5>1$ だから $-2x>\dfrac{3}{2}$
  \boxanswer{$x<-\dfrac{3}{4}$}
\end{answerbox}

\vfill
\nexttime{対数とは何かを学び，指数との関係を考えてみよう！}

\end{document}
